\section{Counting polars of a smooth hypersurface}\label{sec:polar}

We now prove Theorem~\ref{thm:smooth-border}. A smooth degree-$r$
hypersurface in $\PP^{d-1}$ supplies $r(r-1)^{d-2}$ simple intersections
with first polars. These points persist for one sufficiently close affine
determinant. Right and left kernel variables then express each selected
intersection by bilinear equations, whose multidegrees give the numerical
comparison stated in Section~\ref{sec:models}. The determinant kernel-incidence
count and its Schur-complement reduction
build on \cite[Theorem~3 and Section~3.1]{SheshadriBorder2026}. We give the
argument needed here for finitely many simple polar intersections,
including unrestricted affine limits and possible excess components of
the incidence.

For a constant vector $w\in\C^N$, write
\[
 \partial_w F=\sum_{j=1}^N w_j\frac{\partial F}{\partial x_j}.
\]
We call this an ambient first polar: after restricting to an affine slice,
the direction $w$ need not be tangent to the slice. A common zero of $s$
equations in $s$ local coordinates is \emph{simple} if its full Jacobian
is invertible.

\subsection{Simple polar intersections}

The numerical count below is the classical top polar degree of a smooth
hypersurface \cite[Corollary~3.7]{Piene1978}. We give the transversality
argument to retain the individual simple points needed in the limit.
General-transversality principles of this kind are developed in
\cite[Corollary~5]{Kleiman1974}; the argument below works directly over
$\C$ with the incidence of polar directions.

\begin{lemma}\label{polar:simple-points}
Let $r,d\ge2$, and let $G\in\C[x_1,\ldots,x_d]$ be homogeneous of degree
$r$. Suppose that $Y=V(G)\subset\PP^{d-1}$ is a smooth hypersurface.
There are constant directions $w_1,\ldots,w_{d-2}\in\C^d$ for which
\begin{equation}\label{polar:projective-certificate}
 G=0,\qquad \partial_{w_i}G=0\quad(1\le i\le d-2)
\end{equation}
has exactly $\mu=r(r-1)^{d-2}$ distinct simple projective solutions.
There is a linear form $L_0$ nonzero at every one of these points.
\end{lemma}

\begin{proof}
The sections $\partial_wG$ of $\cO_Y(r-1)$ have no common zero on $Y$:
otherwise all partial derivatives of $G$ would vanish at a point of the
smooth hypersurface. We spell out why general choices give simple points.
If $d=2$, the smooth degree-$r$ hypersurface in $\PP^1$ already consists
of $r$ distinct points, and there are no polar equations. Assume $d\ge3$.

Put $\mathcal P=(\C^d)^{d-2}$ and consider
\[
 \mathcal I=\{(x,w_1,\ldots,w_{d-2})\in Y\times\mathcal P:
                  \partial_{w_i}G(x)=0\text{ for every }i\}.
\]
For each $x$, the equation in each separate $w_i$ block is a nonzero
linear equation. Thus $\mathcal I$ is the total space of a vector bundle
of rank $(d-1)(d-2)$ over $Y$. In particular it is smooth, and
$\dim\mathcal I=d(d-2)=\dim\mathcal P$. The incidence is closed in
$Y\times\mathcal P$, so its projection to $\mathcal P$ is proper because
$Y$ is projective. It is surjective: $G$ and any
$d-2$ chosen polars are $d-1$ positive-degree homogeneous equations in
$d$ variables, so the projective dimension theorem gives a common
projective zero.

Characteristic-zero generic smoothness makes the generic fiber of each
dominating component smooth. Remove the proper images of the
non-dominating components and of the locus where the projection is not
smooth. Properness makes these images closed, and generic smoothness
ensures that they omit the generic point of $\mathcal P$. The resulting
nonempty open subset has smooth zero-dimensional fibers. Each such fiber
is a finite set of reduced points. At each point the differentials of the polar
sections on $Y$ are independent. Since $dG\ne0$ in local projective
coordinates, the full system~\eqref{polar:projective-certificate} is
simple there.

This is a proper intersection of hypersurfaces of degrees
$r,r-1,\ldots,r-1$ in $\PP^{d-1}$. Projective B\'ezout therefore gives
exactly $r(r-1)^{d-2}$ points, each of multiplicity one.
Finally, the linear forms vanishing at at least one of these finitely
many points form a finite union of proper hyperplanes in $(\C^d)^*$.
Choose $L_0$ outside that union.
\end{proof}

\subsection{A count that permits excess components}

The kernel equations below need not meet properly everywhere, so we count
only their selected simple zeros. The coefficient in the next lemma is
the multihomogeneous B\'ezout number. Its standard specialization bound
controls nonsingular isolated zeros even when other components are present
\cite[Sections~3--3.2]{Wampler1996}. We give a direct proof of the bilinear
case needed here.

\begin{lemma}\label{polar:product-count}
Let $Y=\prod_{j=1}^s\PP^{a_j}$, where $a_j\ge0$ and
$M=\sum_j a_j\ge1$. For each $1\le\ell\le M$, let $F_\ell$ be a
multihomogeneous polynomial of degree one in two distinct coordinate
groups, indexed by $u_\ell,v_\ell$, and degree zero in the other groups.
Any finite collection of distinct simple common zeros of these equations
has size at most
\begin{equation}\label{polar:product-coefficient}
 [H_1^{a_1}\cdots H_s^{a_s}]
       \prod_{\ell=1}^M(H_{u_\ell}+H_{v_\ell}).
\end{equation}
Other zeros and components of any dimension are allowed.
\end{lemma}

\begin{proof}
First construct a system with the displayed number of zeros, all simple.
Give each occurrence of a coordinate group in an equation a different
label. On $\PP^{a_j}$ choose the labeled linear forms so that every
subfamily of at most $a_j+1$ coefficient vectors is independent.
In coordinates $[z_0:\cdots:z_{a_j}]$, forms
\[
 z_0+t z_1+\cdots+t^{a_j}z_{a_j}
\]
with distinct values of $t$ have this property by the Vandermonde
determinant. When $a_j=0$, use the nowhere-zero form $z_0$ on $\PP^0$.
Replace equation $\ell$ by $Q_\ell=L_\ell R_\ell$, using its two
labeled forms on factors $u_\ell$ and $v_\ell$.

At a common zero of all $Q_\ell$, let $b_j$ count the labeled forms
vanishing in factor $j$. General position implies $b_j\le a_j$.
Each of the $M$ equations requires at least one vanishing form, and no
label occurs in two equations. Hence
\[
 M\le\sum_j b_j\le\sum_j a_j=M.
\]
Exactly one labeled factor vanishes in each equation, and exactly $a_j$
forms vanish in factor $j$. Conversely, choosing one factor of each
equation, with exactly $a_j$ choices assigned to factor $j$, determines
one point of $Y$: the selected $a_j$ hyperplanes in that factor meet in
one point. No additional labeled form vanishes there, because $a_j+1$
such hyperplanes have empty intersection. Different choices therefore
give different points. Expanding the product in
\eqref{polar:product-coefficient} counts exactly these choices; denote
that coefficient by $B$.

Every comparison zero is simple. In equation $\ell$, one factor is zero
and the other is nonzero, so its differential is a nonzero multiple of
the selected linear form's differential. In each projective factor the
selected hyperplanes have independent differentials. The full Jacobian
is therefore invertible. This construction also covers $B=0$, when the
comparison system has no common zero.

It remains to compare an arbitrary system with this one. Interpolate by
\[
 F_{\ell,t}=(1-t)Q_\ell+tF_\ell\qquad(t\in\C).
\]
The common-zero incidence is closed in $Y\times\C$. On it, rank
deficiency of the differential in the $Y$ directions is a closed
algebraic condition. Indeed, it is given by maximal Jacobian minors in
each projective chart. Coordinate changes act invertibly on the
differential, and changing a trivialization multiplies an equation by a
unit. At a common zero the derivative of that unit contributes zero,
so the rank condition is independent of the chart.

Since $Y$ is projective, the parameter values having a rank-deficient
common zero form a Zariski-closed subset $\mathcal B\subset\C$.
The comparison system shows that $0\notin\mathcal B$, so $\mathcal B$
is finite. For $t\notin\mathcal B$, all zeros are isolated by the
implicit function theorem. Their set is closed in the compact space
$Y$, hence finite. The number of zeros is locally constant on
$\C\setminus\mathcal B$: the implicit function theorem preserves each
zero in a disjoint neighborhood, and compactness prevents additional
zeros from appearing outside those neighborhoods as the parameter varies.
The same argument covers an empty fiber. Because the complement of a
finite subset of $\C$ is connected, this number is always the number at
$t=0$, namely $B$.

Finally, choose any finite set of simple zeros at $t=1$. They persist in
disjoint neighborhoods for parameters sufficiently close to $1$. Such
parameters can be chosen outside $\mathcal B$, where the total count is
$B$. The selected set has at most $B$ members. This conclusion does not
assign multiplicities to, or impose conditions on, any excess components
at $t=1$.
\end{proof}

\subsection{Ambient polars of affine determinant limits}

We next apply the product count to an arbitrary affine determinant
approximation. The directions in the statement remain ambient directions;
this permits us to use a projective polar certificate on an affine chart.

\begin{proposition}\label{polar:ambient-bound}
Let $N,n,\mu\ge1$ and $h\ge0$ be integers. Suppose that
$P\in\C[x_1,\ldots,x_N]$ is a coefficientwise limit of determinants of
$n\times n$ matrices of affine-linear forms. Let
$\Psi:\C^{h+1}\to\C^N$ be an affine map and
$w_1,\ldots,w_h\in\C^N$ constant directions. If
\begin{equation}\label{polar:affine-certificate}
 P(\Psi(\xi))=0,\qquad
 (\partial_{w_i}P)(\Psi(\xi))=0\quad(1\le i\le h)
\end{equation}
has at least $\mu$ distinct simple zeros, then
\begin{align}
 \mu&\le [H^{h+1}U^{n-1}V^{n-1}]
            (H+U)^n(H+V)^{n-1}(U+V)^h\label{polar:ambient-coefficient}\\
  &\le 2^h\binom{2n-1}{h+1},\notag\\
 n&\ge\frac{h+1}{4e}\mu^{1/(h+1)}.\label{polar:ambient-size}
\end{align}
The binomial coefficient is zero if $h+1>2n-1$.
\end{proposition}

\begin{proof}
\emph{Choose one sufficiently close determinant.}
Fix $\mu$ distinct simple zeros of~\eqref{polar:affine-certificate}.
All output polynomials lie in the finite-dimensional coefficient space
of degree at most $n$. With $\Psi$ and the $w_i$ fixed, the map
\[
 (Q,\xi)\longmapsto
   \bigl(Q(\Psi(\xi)),(\partial_{w_1}Q)(\Psi(\xi)),\ldots,
                         (\partial_{w_h}Q)(\Psi(\xi))\bigr)
\]
is polynomial in the coefficients of $Q$ and the coordinates of $\xi$.
At every selected zero for $P$, its $\xi$-Jacobian is invertible.
The complex implicit function theorem preserves that simple zero for all
sufficiently close coefficient vectors. Take disjoint neighborhoods of
the $\mu$ zeros and intersect their finitely many coefficient neighborhoods.
There is an actual determinant in the resulting neighborhood,
\begin{equation}\label{polar:chosen-matrix}
 Q(x)=\det A(x),\qquad
 A(x)=A_0+\sum_{j=1}^N x_jA_j,\qquad A_j\in\Mat_n(\C),
\end{equation}
for which the same polar system has $\mu$ distinct simple zeros.
This choice constrains only the coefficients of $Q$. The matrix entries
may be arbitrary complex numbers, however large; no bound on an
approximation order or on coefficients of a parametrized family is used.

\emph{Choose a common chart for the right kernels.}
At each selected zero, the first row of the polar Jacobian is nonzero,
so $d(Q\circ\Psi)\ne0$. If $A(\Psi(\xi))$ had rank at most $n-2$,
its adjugate would vanish and the determinant derivative in every direction
would be zero. Since its determinant is zero, its rank is exactly $n-1$.
It has unique right and left projective kernel lines.

Choose a linear functional nonzero on every one of the finitely many
right kernel lines. Such functionals are the complement of a finite union
of proper hyperplanes. A fixed change of basis in the matrix domain makes
this functional the first coordinate. Equivalently, replace $A$ by $AS$
for a suitable invertible constant matrix $S$. This multiplies $Q$ and
all its polars by the same nonzero constant and preserves simplicity.
Continue to write $A,Q$ for the resulting matrix and polynomial. Every
selected right kernel now has a representative with first coordinate one.
Put
\[
 B_i=\partial_{w_i}A=\sum_{j=1}^N(w_i)_j A_j.
\]
These are constant matrices, and
$\partial_{w_i}Q=\tr(\adj(A)B_i)$.

\emph{Lift to a square system of bilinear equations.}
Let $A^{\mathrm{hom}}(\xi_0,\xi)$ be the entrywise homogenization of the affine matrix
$A\circ\Psi$, so $A^{\mathrm{hom}}(1,\xi)=A(\Psi(\xi))$. On
\[
 \PP^{h+1}_{[\xi_0:\xi]}\times\PP^{n-1}_p\times\PP^{n-1}_b
\]
use the equations
\begin{equation}\label{polar:kernel-system}
 A^{\mathrm{hom}} p=0,\qquad
 (b^{\mathsf T}A^{\mathrm{hom}})_j=0\quad(2\le j\le n),\qquad
 b^{\mathsf T}B_i p=0\quad(1\le i\le h).
\end{equation}
There are $n+(n-1)+h=h+2n-1$ equations, equal to the product dimension.
In the chart $p_1\ne0$, the omitted first left-kernel equation follows
from the others, since
\[
 0=b^{\mathsf T}(A^{\mathrm{hom}} p)
   =p_1(b^{\mathsf T}A^{\mathrm{hom}})_1+
       \sum_{j=2}^n p_j(b^{\mathsf T}A^{\mathrm{hom}})_j.
\]
At a selected point, rank $n-1$ implies
$\adj(A)=\rho p b^{\mathsf T}$ for normalized kernel representatives
and some $\rho\ne0$: every column of the adjugate is in the right kernel,
every row is in the left kernel, and some cofactor is nonzero.
Consequently $\partial_{w_i}Q=\rho b^{\mathsf T}B_i p$ there.
Each selected affine zero therefore has a unique lift satisfying
\eqref{polar:kernel-system}.

The lift is not yet known to be simple. We prove this locally before
applying Lemma~\ref{polar:product-count}; other components of the kernel
system will impose no further requirement.

\begin{samepage}
\emph{Eliminate the kernel coordinates.}
Fix one selected lift and normalize $p_1=1$. Choose a nonzero coordinate
of its left kernel vector, move that row to the first position, and
normalize the corresponding left coordinate to one. These are local
coordinate choices, which can differ between selected points. The
cofactor obtained by deleting this first row and first column is nonzero,
because it is the corresponding entry of $\rho p b^{\mathsf T}$.
In these row orders write the ambient matrix near $\Psi(\xi)$ as
\[
 A(x)=\begin{pmatrix}a(x)&\beta(x)\\\gamma(x)&E(x)\end{pmatrix},
 \qquad
 p=\binom{1}{u},\qquad b=\binom{1}{v},
\]
where $E$ has size $(n-1)\times(n-1)$. The chosen cofactor makes $E$
invertible at the selected point, so restrict to an ambient neighborhood
on which it is invertible. When $n=1$, $E$ is the empty matrix with
determinant one, and the auxiliary vectors $u,v$ are empty.

\end{samepage}
The same row permutation is applied to every $B_i$. The last $n-1$ right
equations and all retained left equations give
\[
 \gamma+Eu=0,\qquad \beta+v^{\mathsf T}E=0.
\]
Their auxiliary Jacobian has invertible diagonal blocks $E,E^{\mathsf T}$,
so they eliminate the kernel coordinates holomorphically as
\begin{equation}\label{polar:solved-kernels}
 u=-E^{-1}\gamma,\qquad v^{\mathsf T}=-\beta E^{-1}.
\end{equation}
The remaining right equation, evaluated on the slice, is
\[
 s(\xi)=S(\Psi(\xi))=0,
 \qquad S(x)=a(x)-\beta(x)E(x)^{-1}\gamma(x).
\]
The Schur determinant identity gives
\begin{equation}\label{polar:schur-determinant}
 Q(x)=\sigma\det(E(x))S(x),
\end{equation}
where $\sigma\in\{1,-1\}$ is the sign of the row permutation. It follows
by subtracting $\beta E^{-1}$ times the lower block rows from the first
row. Here $S$ is defined on an ambient neighborhood where $E$ is
invertible, so it can be differentiated in each ambient direction $w_i$.

After substituting~\eqref{polar:solved-kernels}, let
$g_i(\xi)$ denote the bilinear equation $b^{\mathsf T}B_i p$.
For any constant ambient direction $w$, differentiation of $S$ gives
\begin{align*}
 \partial_w S
 &=\partial_w a-(\partial_w\beta)E^{-1}\gamma
       +\beta E^{-1}(\partial_w E)E^{-1}\gamma
       -\beta E^{-1}(\partial_w\gamma)\\
 &=(1,-\beta E^{-1})(\partial_w A)
                 \binom{1}{-E^{-1}\gamma}.
\end{align*}
Thus $g_i=(\partial_{w_i}S)\circ\Psi$. In particular this computation
does not identify an ambient derivative with a derivative along the slice.
If $f=Q\circ\Psi$ and $q_i=(\partial_{w_i}Q)\circ\Psi$, differentiating
\eqref{polar:schur-determinant} in the ambient space yields
\begin{equation}\label{polar:reduced-system}
 \begin{gathered}
 f=c\,s,\qquad q_i=c\,g_i+c_i s,\\
 c=\sigma(\det E)\circ\Psi\ne0,\qquad
 c_i=\sigma(\partial_{w_i}\det E)\circ\Psi.
 \end{gathered}
\end{equation}
The original equation vector $(f,q_1,\ldots,q_h)$ and the reduced vector
$(s,g_1,\ldots,g_h)$ therefore differ by a triangular holomorphic matrix
with diagonal entries $c$. At a common zero their Jacobians differ by
the value of that invertible matrix. The original zero was simple, so
the reduced Jacobian is invertible.

Finally, a tangent vector annihilated by the full kernel-system Jacobian
first satisfies the differentials of the eliminated equations. Their
invertible auxiliary block determines its $(u,v)$ component from its
$\xi$ component. The remaining equations and the invertible reduced
Jacobian force that $\xi$ component to be zero, and hence the auxiliary
component is zero as well. The full square Jacobian is invertible.
All $\mu$ selected lifts are thus distinct simple zeros of
\eqref{polar:kernel-system}.

\emph{Count the selected lifts.}
The three groups of equations in~\eqref{polar:kernel-system} have
multidegrees and multiplicities
\[
 \begin{array}{c|c|c}
 \text{equations}&\text{multidegree}&\text{number}\\ \hline
 A^{\mathrm{hom}} p &(1,1,0)&n\\
 (b^{\mathsf T}A^{\mathrm{hom}})_j, j\ge2 &(1,0,1)&n-1\\
 b^{\mathsf T}B_i p &(0,1,1)&h.
 \end{array}
\]
Lemma~\ref{polar:product-count} gives the first inequality
in~\eqref{polar:ambient-coefficient}, regardless of other zeros or
components. The product has nonnegative coefficients. Setting $U=V=1$
sums the coefficients with the same $H$-exponent, and hence bounds the
particular coefficient by
\[
 [H^{h+1}]\,2^h(H+1)^{2n-1}
       =2^h\binom{2n-1}{h+1}.
\]
Since $\mu>0$, we must have $h+1\le2n-1$. For integers $1\le b\le a$,
\[
 \binom ab\le\frac{a^b}{b!}\le\left(\frac{ea}{b}\right)^b;
\]
the last bound follows from
$\log(b!)\ge\int_1^b\log t\,dt\ge b\log b-b$.
Taking $(h+1)$st roots gives
\[
 \mu^{1/(h+1)}
 \le2^{h/(h+1)}\frac{e(2n-1)}{h+1}
 \le\frac{4en}{h+1},
\]
which is~\eqref{polar:ambient-size}. No comparison between $n$ and
the number $N$ of ambient variables, or between $n$ and $h+2$, was assumed.
\end{proof}

\subsection{The smooth-hypersurface lower bound}

We have proved that a fixed finite set of simple polar intersections
persists under arbitrary affine determinant approximation, and that its
kernel lifts satisfy the multidegree bound. We now use the certificate
provided by the smooth hypersurface.

\begin{proof}[Proof of Theorem~\ref{thm:smooth-border}]
Let $n$ be any size for which $G$ is a coefficientwise limit of affine
$n\times n$ determinants. Choose the directions and $L_0$ from
Lemma~\ref{polar:simple-points}. The affine hyperplane $L_0=1$ identifies
with the projective chart $L_0\ne0$ and contains representatives of all
$\mu=r(r-1)^{d-2}$ selected points. Parametrize it by an affine isomorphism
\[
 \Psi:\C^{d-1}\longrightarrow\{x\in\C^d:L_0(x)=1\}.
\]
Restricting the homogeneous equations to this hyperplane is precisely
their dehomogenization in that projective chart. Thus the equations
\[
 G\circ\Psi=0,\qquad
 (\partial_{w_i}G)\circ\Psi=0\quad(1\le i\le d-2)
\]
have $\mu$ distinct simple affine solutions. The directions remain constant
vectors in $\C^d$; they need not be tangent to $L_0=1$.
Apply Proposition~\ref{polar:ambient-bound} with $N=d$, $h=d-2$, and
$P=G$. It gives
\[
 n\ge\frac{d-1}{4e}
       \bigl(r(r-1)^{d-2}\bigr)^{1/(d-1)}
   \ge\frac{(r-1)(d-1)}{4e},
\]
where the last inequality uses $r\ge r-1>0$. This holds for every
admissible size $n$, so taking the minimum proves the theorem.
\end{proof}
